\documentclass[twocolumn,11pt,a4paper]{article}

%% -----------------------------------------------------------------------
%%  Packages
%% -----------------------------------------------------------------------
\usepackage{amsmath,amssymb,amsthm}
\usepackage{mathtools}
\usepackage{lmodern}
\usepackage[protrusion=true,expansion=false]{microtype}
\usepackage{geometry}
\geometry{a4paper, top=2.2cm, bottom=2.2cm, left=1.8cm, right=1.8cm, columnsep=0.7cm}
\usepackage{hyperref}
\hypersetup{colorlinks=true,linkcolor=blue,citecolor=blue,urlcolor=blue}
\usepackage{cleveref}
\usepackage{booktabs}
\usepackage{array}
\usepackage{enumitem}
\usepackage{xcolor}
\usepackage{natbib}
\bibpunct{[}{]}{,}{n}{}{;}
\usepackage{titlesec}
\usepackage{abstract}
\usepackage{graphicx}

%% -----------------------------------------------------------------------
%%  Theorem environments
%% -----------------------------------------------------------------------
\newtheorem{theorem}{Theorem}[section]
\newtheorem{lemma}[theorem]{Lemma}
\newtheorem{proposition}[theorem]{Proposition}
\newtheorem{corollary}[theorem]{Corollary}
\newtheorem{definition}[theorem]{Definition}
\newtheorem{construction}[theorem]{Construction}
\newtheorem{example}[theorem]{Example}
\theoremstyle{remark}
\newtheorem{remark}[theorem]{Remark}
\newtheorem{econinterp}[theorem]{Economic Interpretation}
\newtheorem{comparison}[theorem]{Comparison}

%% -----------------------------------------------------------------------
%%  Mathematical macros (inherited)
%% -----------------------------------------------------------------------
\newcommand{\R}{\mathbb{R}}
\newcommand{\N}{\mathbb{N}}
\newcommand{\Sflat}{S_{\flat}}
\newcommand{\Sig}{\Sigma}
\newcommand{\Rec}{\mathcal{R}}
\newcommand{\Meth}{\mathcal{M}}
\newcommand{\Act}{\mathrm{Act}}
\newcommand{\Reach}{\mathrm{Reach}}
\newcommand{\tol}{\tau}
\newcommand{\Ens}{\mathcal{E}}
\newcommand{\Compose}{\diamond}

%% -----------------------------------------------------------------------
%%  New macros (this paper)
%% -----------------------------------------------------------------------
\newcommand{\Pric}{P}                    % price functional
\newcommand{\MidP}{p^*}                  % mid-market price
\newcommand{\Bid}{B^*}                   % bid price
\newcommand{\Ask}{A^*}                   % ask price
\newcommand{\SpreadT}{\Delta^T}          % trading spread
\newcommand{\SpreadI}{\Delta^I}          % information spread
\newcommand{\ValT}{V^T}                  % trading value
\newcommand{\ValI}{V^I}                  % informational value
\newcommand{\FV}{V^*}                    % fundamental value
\newcommand{\SflatT}{S_{\flat,\mathrm{agg}}}        % aggregate (min) floor
\newcommand{\SflatI}{S_{\flat,\mathrm{comp}}}       % composite (product) floor
\newcommand{\SflatTinf}{S_{\flat,\mathrm{agg},\infty}}  % asymptotic aggregate floor
\newcommand{\Prem}{\Pi}                  % information premium
\newcommand{\logfl}{\ell}               % log-floor
\newcommand{\srad}{\rho}               % spectral radius (from Paper 3)
\newcommand{\CV}{\mathrm{CV}}          % coefficient of variation (from Paper 3)

%% -----------------------------------------------------------------------
%%  Title block
%% -----------------------------------------------------------------------
\title{\textbf{Mathematical Foundation for Price:\\
Cell-Values, Dual Spreads, and Fundamental Value\\
in the Receiver-Agent Framework}}

\author{Kundai Farai Sachikonye\\[4pt]
\small Technical University of Munich\\
\small AIMe Registry for Artificial Intelligence\\
\small Maximus-von-Imhof-Forum 3, 85354 Freising, Germany\\
\small \texttt{kundai.sachikonye@wzw.tum.de}}

\date{}

\begin{document}

\maketitle

%% -----------------------------------------------------------------------
%%  Abstract
%% -----------------------------------------------------------------------
\begin{abstract}
\noindent
This paper establishes a rigorous mathematical foundation for price in the
receiver-agent framework of the companion
papers~\cite{sachikonye2025agents,sachikonye2025market,sachikonye2025het}.
We derive price as a theorem rather than an assumption, proving ten principal
results:
(i)~the \emph{Price-Cell Theorem} --- every price is a cell $C(p,\tol)$ in
outcome space; point prices require $\beta=0$, which the Floor Theorem of
Paper~1 forbids for any bounded receiver;
(ii)~the \emph{Cell-Value Decomposition} --- each agent $A$ assigns two
values to every cell $C$: a \emph{trading value}
$\ValT_A(C) = \tol(C) - \SflatT(\Ens)$ (set by the best agent's floor) and
an \emph{informational value}
$\ValI_A(C) = \tol(C) - \SflatI(\Ens)$ (set by the ensemble's composite
floor);
(iii)~the \emph{Dual-Spread Theorem} --- the trading spread $\SpreadT = 2
\SflatT(\Ens) = 2\min_i\beta_i$ and the information spread
$\SpreadI = 2\SflatI(\Ens)$ are two distinct measures of market quality,
with $\SpreadI \le \SpreadT$ for $n \ge 2$;
(iv)~the \emph{Two-Value Theorem} (main result) --- the informational value
always exceeds the trading value: $\ValI_\Ens(C) \ge \ValT_\Ens(C)$,
with the gap $\Prem = \ValI - \ValT = \SflatT(\Ens) - \SflatI(\Ens) > 0$
the \emph{information premium} of the ensemble;
(v)~the \emph{Equilibrium Price Theorem} --- the equilibrium price $\MidP$
is the center of the purpose cell $C^*$, unique by the Banach contraction
of Paper~2;
(vi)~the \emph{Price Discovery Rate Theorem} --- price converges to $\MidP$
at rate $\srad_\Ens = \SflatT(\Ens)/\Sigma$, matching the spectral radius
of Paper~3;
(vii)~the \emph{Fundamental Value Theorem} --- for a Class~$\mathcal{C}_E$
(Borel--Cantelli efficient) ensemble, the sequence of equilibrium prices
$\MidP_n$ converges to the \emph{fundamental value} $\FV$, the unique
point at which perfect price discovery would occur;
(viii)~the \emph{No-Arbitrage Theorem} --- a market equilibrium exists if
and only if there is a cell $C$ with $\tol(C) > \SflatT(\Ens)$; failure
of this condition produces a ``no-trade zone'' with no equilibrium;
(ix)~the \emph{Transaction Theorem} --- buyer and seller transact if and
only if their reachability sets for the equilibrium cell $C^*$ overlap;
the transaction interval is $[\MidP - \ValT_{\mathrm{seller}},
\MidP + \ValT_{\mathrm{buyer}}]$;
and
(x)~the \emph{Law of One Price Theorem} --- for agents sharing the same
knowledge framework $K$, simultaneous attainment of the same cell
$C$ implies agreement on the decoded value: $D_i(x) = D_j(x)$ for all
$x \in C$.
Forty-five numerical experiments verify all claims to machine precision.
\end{abstract}

\tableofcontents

\bigskip

%% ======================================================================
\section{Introduction}
\label{sec:intro}
%% ======================================================================

\subsection{The Price Question}
\label{sec:intro:price-q}

Price is the central object of economics.  Yet the theoretical foundations
of price remain contested.  Classical models define price as a point
$p^* \in \R$, derived either as a Walrasian market-clearing price
vector~\cite{arrowdebreu1954,debreu1959}, a Nash equilibrium payoff~\cite{nash1950},
or a Black--Scholes fair value~\cite{blackscholes1973}.  All three
approaches assume that agents can specify and respond to point prices with
arbitrary precision — an assumption that, in the language of the companion
papers, requires each agent's receiver floor $\beta$ to be exactly zero.

The Floor Theorem of Paper~1~\cite{sachikonye2025agents} proves that
$\beta > 0$ for every bounded receiver.  Point prices are therefore
unattainable by any bounded economic agent.  This is not a technical
inconvenience but a structural feature: the positivity of the floor is the
mathematical fingerprint of bounded rationality.

Papers~2 and~3~\cite{sachikonye2025market,sachikonye2025het} established
that ensembles of bounded agents nevertheless reach equilibrium — the
purpose cell $C^*$ — and that the information efficiency of this equilibrium
is governed by two distinct floor quantities: the \emph{aggregate floor}
$\SflatT(\Ens) = \min_i \Sflat(A_i)$ and the \emph{composite floor}
$\SflatI(\Ens) = \prod_i \Sflat(A_i)/\Sigma^{n-1}$.

The present paper asks: given that equilibrium is a cell rather than a
point, what is ``the price''?  How do the two floor quantities translate
into observable market phenomena?  And what is the relationship between the
equilibrium cell and the classical notion of a fundamental value?

\subsection{Main Contributions}
\label{sec:intro:contrib}

\begin{enumerate}[leftmargin=*, label=\textbf{(\roman*)}]
\item \textbf{Price-Cell Theorem}
  (\cref{thm:price-cell}).
  Price is not a primitive but a derived object: a cell $C(p,\tol)$ in
  outcome space $X$.  Point prices are unattainable for any bounded agent.

\item \textbf{Cell-Value Decomposition}
  (\cref{thm:cell-value}).
  Agent $A$ assigns two values to cell $C$:
  $\ValT_A(C) = \tol(C) - \Sflat(A)$ (trading value, positive iff the
  agent can attain $C$) and $\ValI_\Ens(C) = \tol(C) - \SflatI(\Ens)$
  (informational value, set by the ensemble's composite floor).

\item \textbf{Dual-Spread Theorem}
  (\cref{thm:dual-spread}).
  Trading spread $\SpreadT = 2\min_i\Sflat(A_i)$ is set by the best agent
  alone.  Information spread $\SpreadI = 2\SflatI(\Ens)$ is set by all
  agents jointly.  $\SpreadI \le \SpreadT$ for $n \ge 2$, with equality
  iff $n=1$.

\item \textbf{Two-Value Theorem} (Main result)
  (\cref{thm:two-value}).
  $\ValI_\Ens(C) \ge \ValT_\Ens(C)$ for every cell $C$ and every ensemble
  $\Ens$.  The \emph{information premium} $\Prem = \ValI - \ValT =
  \SflatT - \SflatI \ge 0$ is zero iff $n=1$ and strictly positive for
  all $n\ge 2$.

\item \textbf{Equilibrium Price Theorem}
  (\cref{thm:equil-price}).
  The equilibrium price $\MidP$ is the center of the purpose cell $C^* =
  B(\MidP, r^*)$.  $\MidP$ is unique given $C^*$, which is unique by the
  Banach contraction of Paper~2.

\item \textbf{Price Discovery Rate Theorem}
  (\cref{thm:discovery-rate}).
  Starting from any initial price $p_0$, the price discovery sequence
  $p_t \to \MidP$ at rate $\srad_\Ens = \SflatT(\Ens)/\Sigma$, matching
  the spectral radius of the purpose-cell contraction (Paper~3).

\item \textbf{Fundamental Value Theorem}
  (\cref{thm:fv}).
  For a Borel--Cantelli-efficient ensemble ($\Ens_\infty \in \mathcal{C}_E$),
  $\SflatI(\Ens_n) \to 0$ and the equilibrium cell collapses to a point:
  $C^*_n \to \{\FV\}$.  The fundamental value $\FV$ is the unique limit.

\item \textbf{No-Arbitrage Theorem}
  (\cref{thm:no-arb}).
  A purpose cell exists iff $\SflatT(\Ens) < \Sigma$.  When this fails,
  no equilibrium price exists and the market is in a \emph{no-trade zone}.
  No-arbitrage is equivalent to positive purpose margin.

\item \textbf{Transaction Theorem}
  (\cref{thm:transaction}).
  Buyer $A_B$ and seller $A_S$ transact at equilibrium cell $C^*$ iff
  $\Sflat(A_B) + \Sflat(A_S) \le 2\tol(C^*)$.  The transaction interval
  is non-empty iff this condition holds.

\item \textbf{Law of One Price Theorem}
  (\cref{thm:loop}).
  Agents with the same knowledge framework $K_i = K_j$ who simultaneously
  attain cell $C$ must agree on their decoded representation at every state
  in $C$: $D_i(x) = D_j(x)$ for all $x \in C$.  Distinct prices within
  the cell do not violate the law of one price; they reflect floor-bounded
  measurement noise.
\end{enumerate}

\subsection{Organization}
\label{sec:intro:org}

\Cref{sec:prelim} recaps the inherited framework.
\Cref{sec:price-cells} introduces price cells and proves the Price-Cell
Theorem.
\Cref{sec:cell-value} develops cell-value functions and the Cell-Value
Decomposition.
\Cref{sec:dual-spread} proves the Dual-Spread Theorem and Two-Value Theorem.
\Cref{sec:equil-price} establishes the Equilibrium Price Theorem.
\Cref{sec:discovery} proves the Price Discovery Rate Theorem.
\Cref{sec:fv} develops the Fundamental Value Theorem.
\Cref{sec:no-arb} proves the No-Arbitrage Theorem.
\Cref{sec:transaction} proves the Transaction Theorem.
\Cref{sec:loop} proves the Law of One Price Theorem.
\Cref{sec:numerical} presents forty-five numerical experiments.
\Cref{sec:literature} surveys connections to the literature.
\Cref{sec:conclusion} states conclusions.

%% ======================================================================
\section{Inherited Framework}
\label{sec:prelim}
%% ======================================================================

We recall the primitives from the companion
papers~\cite{sachikonye2025agents,sachikonye2025market,sachikonye2025het}.
The \emph{outcome space} $(X, d, \Sigma)$ is a bounded metric space with
maximal semantic distance $\Sigma$.  A \emph{receiver} $\Rec = (K, D, \Pi,
\beta)$ has decoder $D: X \to K$, projection $\Pi: K \to 2^X$, and floor
$\beta > 0$.  The S-functional is $S(\Rec, x; C) = \inf_{c \in \Pi(D(x))}
d(c, C) + \beta$ (reduced to $\max(\beta, d(x, C))$ for ball cells with ball
projections).

An agent triple $A = (\Rec, \Meth, \mathcal{G})$ has floor
$\Sflat(A) = \beta \cdot \Sflat(\Meth)/\Sigma$.  An ensemble $\Ens_n =
(A_1,\ldots,A_n)$ has:
\begin{align}
  \SflatT(\Ens_n) &:= \min_{1\le i\le n} \Sflat(A_i)
  && \text{(aggregate floor)}, \label{eq:agg}\\
  \SflatI(\Ens_n) &:= \frac{\prod_{i=1}^n \Sflat(A_i)}{\Sigma^{n-1}}
  && \text{(composite floor)}. \label{eq:comp}
\end{align}
The multi-agent S-functional $S(\Ens_n, x; C) = \min_i S(A_i, x; C)$ has
floor $\SflatT(\Ens_n)$.  The purpose cell $C^*$ from Paper~2 satisfies
$\Phi_{\Ens_n}(C^*) = \SflatT(\Ens_n)$ and $\tol(C^*) > \SflatT(\Ens_n)$,
proved by Banach contraction on $(\mathrm{Cell}(X), d_H)$.  Throughout,
$\Sigma = 100$ is the canonical constant.

\begin{remark}[Two floors]
The aggregate floor $\SflatT$ governs the ensemble's ability to attain a
cell in a single step (one agent fires).  The composite floor $\SflatI$
governs the ensemble's collective information about the cell (all agents
contribute).  This distinction, first noted in
Paper~2~\cite{sachikonye2025market}, is the mathematical origin of the
dual-spread structure proved in this paper.
\end{remark}

%% ======================================================================
\section{Price Cells}
\label{sec:price-cells}
%% ======================================================================

\begin{definition}[Price cell]
\label{def:price-cell}
For a price level $p \in X$ and tolerance $\tol > 0$, the
\emph{price cell} is the ball
\[
  C(p, \tol) \;:=\; B(p, \tol) \;=\; \{x \in X : d(x, p) \le \tol\}.
\]
The price level $p$ is the \emph{mid-price} of $C(p,\tol)$ and $\tol$ is
the \emph{half-spread}.  A \emph{point price} is the degenerate case
$C(p, 0) = \{p\}$.
\end{definition}

\begin{theorem}[Price-Cell]
\label{thm:price-cell}
For any bounded receiver $\Rec = (K, D, \Pi, \beta)$ with $\beta > 0$ and
any target cell $C \subseteq X$:
\begin{enumerate}[label=\textup{(\alph*)}]
\item If $\tol(C) = 0$ (point price), then $S(\Rec, x; C) \ge \beta > 0 =
  \tol(C)$ for all $x \in X$.  No bounded agent can ever attain a point price.
\item If $\tol(C) > \beta$, then there exist states $x \in X$ from which
  the agent can attain $C$: $S(\Rec, x; C) \le \tol(C)$.
\item The minimum achievable price tolerance for agent $A$ is $\tol_{\min}(A)
  = \Sflat(A) > 0$.  Any cell with $\tol(C) \le \Sflat(A)$ is unattainable
  by $A$.
\end{enumerate}
\end{theorem}
\begin{proof}
\textit{(a)}  For the ball receiver, $S(\Rec, x; C) = \max(\beta, d(x,C))$.
When $C = \{p\}$, $d(x, C) = d(x, p) \ge 0$, so $S \ge \beta > 0$.
But $\tol(\{p\}) = 0$, so $S(\Rec, x; \{p\}) \ge \beta > 0 = \tol(\{p\})$
for all $x$: the agent never attains $\{p\}$.

\textit{(b)}  For any $x \in C$, $d(x,C) = 0$, so $S(\Rec, x; C) =
\max(\beta, 0) = \beta$.  If $\tol(C) > \beta$, then $S(\Rec, x; C) = \beta
< \tol(C)$: the agent attains $C$ from any in-cell state.

\textit{(c)}  By the Floor Theorem of Paper~1, $\Sflat(A) = \beta\cdot
\Sflat(\Meth)/\Sigma > 0$.  For any cell with $\tol(C) \le \Sflat(A) \le
S(\Rec, x; C)$ for all $x$, the agent cannot attain $C$.
\end{proof}

\begin{remark}[Price uncertainty is irreducible]
Part~(a) of \cref{thm:price-cell} gives the fundamental theorem of price
theory in the receiver-agent framework: price uncertainty $\tol > 0$ is
irreducible for any bounded agent.  The half-spread $\tol$ is not a trading
cost to be minimised but a structural feature of any bounded receiver.
\end{remark}

\begin{figure*}[!htbp]
    \centering
    \includegraphics[width=\textwidth]{panels/paper4_panel_1.png}
    \caption{\textbf{Price Cells and Cell-Value Decomposition: Dual-Value Surface, Trading vs.\ Informational Value, Information Premium Dynamics, and Floor Sensitivity.}
    (\textbf{A})~Three-dimensional surfaces of the trading value $V^T(\mathcal{C}) = \tau - \bar{S}_{\mathrm{agg}}(\mathcal{E})$ (blue) and informational value $V^I(\mathcal{C}) = \tau - \bar{S}_{\mathrm{comp}}(\mathcal{E})$ (orange) over cell radius $\tau \in [5, 60]$ and ensemble size $n \in \{1, \ldots, 11\}$, for a uniform ensemble with $\beta = 15$ and $\Sigma = 100$. 
    (\textbf{B})~Trading value $V^T$ (solid) and informational value $V^I$ (dashed) as functions of cell radius $\tau$ for a three-agent ensemble with floors $\{10, 20, 35\}$.  Both values are non-negative for $\tau$ above the respective floor thresholds: $V^T > 0$ only when $\tau > \bar{S}_{\mathrm{agg}} = 10$, and $V^I > 0$ when $\tau > \bar{S}_{\mathrm{comp}} \approx 0.7$. 
    (\textbf{C})~Information premium $\Pi = \bar{S}_{\mathrm{agg}} - \bar{S}_{\mathrm{comp}}$ as a function of ensemble size $n$ for three floor configurations: uniform $\beta = 15$, linearly increasing $\beta \in [5, 40]$, and decreasing $\beta \in [20, 2]$.   
    (\textbf{D})~Trading value $V^T = \tau - \beta$ for a single agent (solid, $n = 1$, where composite equals aggregate) and informational value $V^I = \tau - \bar{S}_{\mathrm{comp}}(\mathcal{E}_3)$ for a uniform three-agent ensemble (dashed), as functions of agent floor $\beta$ with fixed $\tau = 30$. }
    \label{fig:price_panel_1}
\end{figure*}

\begin{econinterp}
\label{ei:price-cell}
In equity markets, a ``price'' is reported to two decimal places (a tick
of 0.01).  This is precisely a price cell $C(p, 0.005)$ — a ball of radius
half a tick around the stated price.  Any quote more precise than one tick
is operationally meaningless for agents whose floors $\beta_i > 0.005$.
The receiver-agent framework predicts that minimum tick size is endogenous:
it is determined by the floors of the market participants, not by regulatory
fiat.
\end{econinterp}

%% ======================================================================
\section{Cell-Value Functions}
\label{sec:cell-value}
%% ======================================================================

\begin{definition}[Cell-value functions]
\label{def:cell-value}
For an ensemble $\Ens = (A_1,\ldots,A_n)$ and a price cell
$C(p,\tol)$, define:
\begin{align}
  \ValT_\Ens(C) &:= \tol(C) - \SflatT(\Ens) = \tol(C) - \min_i \Sflat(A_i),
  \label{eq:vt}\\
  \ValI_\Ens(C) &:= \tol(C) - \SflatI(\Ens) = \tol(C)
  - \frac{\prod_i \Sflat(A_i)}{\Sigma^{n-1}}.
  \label{eq:vi}
\end{align}
$\ValT_\Ens(C)$ is the \emph{trading value} of $C$ to $\Ens$: the excess
of the cell tolerance over the best agent's floor.
$\ValI_\Ens(C)$ is the \emph{informational value}: the excess over the
composite floor.  Both are positive iff the ensemble can attain $C$.
\end{definition}

\begin{theorem}[Cell-Value Decomposition]
\label{thm:cell-value}
For any ensemble $\Ens_n$ and cell $C$:
\begin{enumerate}[label=\textup{(\alph*)}]
\item $\ValT_\Ens(C) = \max_i [\tol(C) - \Sflat(A_i)]$: trading value is
  the maximum over all agents of the individual purpose margin.
\item $\ValT_\Ens(C) \ge \ValT_{A_i}(C) := \tol(C) - \Sflat(A_i)$ for all $i$.
\item $\ValT_\Ens(C)$ and $\ValI_\Ens(C)$ are both increasing in $\tol(C)$
  and decreasing in the floors $\Sflat(A_i)$.
\item For any single-agent ensemble ($n=1$): $\ValT_A(C) = \ValI_A(C) =
  \tol(C) - \Sflat(A)$.
\end{enumerate}
\end{theorem}
\begin{proof}
\textit{(a)}  $\SflatT(\Ens) = \min_i \Sflat(A_i)$, so $\ValT_\Ens(C) =
\tol - \min_i \Sflat(A_i) = \max_i [\tol - \Sflat(A_i)]$.

\textit{(b)}  $\ValT_\Ens(C) = \max_i[\tol - \Sflat(A_i)] \ge \tol -
\Sflat(A_i) = \ValT_{A_i}(C)$ for each $i$.

\textit{(c)}  Both \eqref{eq:vt} and \eqref{eq:vi} are increasing affine
in $\tol(C)$ and decreasing in each $\Sflat(A_i)$ (since $\SflatT$ is
decreasing in each $\Sflat(A_i)$ and $\SflatI$ is a product).

\textit{(d)}  For $n=1$: $\SflatT(A) = \SflatI(A) = \Sflat(A)$, so both
values coincide.
\end{proof}

%% ======================================================================
\section{Dual Spreads and the Two-Value Theorem}
\label{sec:dual-spread}
%% ======================================================================

\begin{definition}[Trading and information spreads]
\label{def:spreads}
For the ensemble $\Ens_n$ at equilibrium cell $C^* = B(\MidP, r^*)$:
\begin{align}
  \SpreadT &:= 2\SflatT(\Ens_n) = 2\min_i \Sflat(A_i), \label{eq:st}\\
  \SpreadI &:= 2\SflatI(\Ens_n) = 2\frac{\prod_i \Sflat(A_i)}{\Sigma^{n-1}}.
  \label{eq:si}
\end{align}
The \emph{information premium} is $\Prem := \ValI_\Ens(C) - \ValT_\Ens(C)
= \SflatT(\Ens) - \SflatI(\Ens)$ (= half the spread difference).
\end{definition}

\begin{theorem}[Dual-Spread]
\label{thm:dual-spread}
For any $n$-agent ensemble $\Ens_n$ with $n \ge 2$:
\begin{enumerate}[label=\textup{(\alph*)}]
\item $\SpreadI \le \SpreadT$, with equality iff $n=1$.
\item $\SpreadT$ depends only on the single best agent ($\min_i \Sflat(A_i)$).
  Adding any agent with floor above $\min_i \Sflat(A_i)$ does not change
  $\SpreadT$.
\item $\SpreadI$ depends on all agents.  Adding any agent with
  $\Sflat(A_{n+1}) < \Sigma$ strictly reduces $\SpreadI$.
\item The ratio $\SpreadI/\SpreadT = \SflatI(\Ens)/\SflatT(\Ens)
  = \prod_{i} q_i / \min_i q_i \in (0,1]$, where $q_i = \Sflat(A_i)/\Sigma$.
\end{enumerate}
\end{theorem}
\begin{proof}
\textit{(a)}  By AM-GM, for any positive reals $a_1,\ldots,a_n$:
$(\prod_i a_i)^{1/n} \le \frac{1}{n}\sum_i a_i$.  In particular,
$\prod_i a_i \le a_j^{n-1} \cdot a_j = a_j^n$ for $a_j = \min_i a_i$...
more directly: $\SflatI(\Ens) = \prod_i f_i/\Sigma^{n-1}$ where
$f_i = \Sflat(A_i)$ and $\SflatT(\Ens) = \min_i f_i$.
$\SflatI/\SflatT = \prod_i f_i / (\Sigma^{n-1} \min_i f_i)$.
Let $f_* = \min_i f_i$; then $f_i \le \Sigma$ for all $i$, so
$\prod_i f_i \le f_* \cdot \Sigma^{n-1}$, giving
$\SflatI \le f_* \cdot \Sigma^{n-1}/\Sigma^{n-1} = f_* = \SflatT$.
Equality: if $n=1$, $\SflatI = f_1 = \SflatT$.  For $n \ge 2$, equality
requires $\prod_i f_i = f_* \Sigma^{n-1}$, i.e.\ $f_i = \Sigma$ for all
$i \ne \arg\min$, but $f_i < \Sigma$ for all $i$, giving strict inequality.

\textit{(b)--(c)}  $\SflatT = \min_i f_i$: changes only when $f_{n+1} <
\min_i f_i$.  $\SflatI = \prod_i f_i/\Sigma^{n-1}$: multiplied by
$f_{n+1}/\Sigma < 1$ when a new agent joins, so always decreases.

\textit{(d)}  Direct computation.
\end{proof}

\begin{theorem}[Two-Value]
\label{thm:two-value}
For any ensemble $\Ens_n$ and any cell $C$ with $\tol(C) > \SflatT(\Ens_n)$:
\begin{enumerate}[label=\textup{(\alph*)}]
\item $\ValI_\Ens(C) \ge \ValT_\Ens(C) \ge 0$.
\item The information premium $\Prem = \ValI_\Ens(C) - \ValT_\Ens(C) =
  \SflatT(\Ens_n) - \SflatI(\Ens_n) \ge 0$, with equality iff $n=1$.
\item $\Prem$ is strictly increasing in ensemble size $n$ and strictly
  increasing in floor heterogeneity (in the majorization sense of Paper~3).
\item $\Prem/\SflatT(\Ens_n) = 1 - \SflatI(\Ens_n)/\SflatT(\Ens_n) \in [0,1)$
  is the \emph{relative information premium} — the fraction by which the
  composite floor falls below the aggregate floor.
\end{enumerate}
\end{theorem}
\begin{proof}
\textit{(a)}  $\ValI - \ValT = \SflatT - \SflatI \ge 0$ by \cref{thm:dual-spread}(a).
$\ValT \ge 0$ by hypothesis.

\textit{(b)}  Follows from \cref{thm:dual-spread}(a) with equality criterion.

\textit{(c)}  By \cref{thm:dual-spread}(c), $\SflatI$ strictly decreases
as $n$ grows.  By the Heterogeneity Dominance Theorem of Paper~3, $\SflatI$
strictly decreases with heterogeneity.  Since $\SflatT = \min_i f_i$ is
fixed once the best agent is included, $\Prem = \SflatT - \SflatI$ strictly
increases.

\textit{(d)}  $\Prem/\SflatT = (\SflatT - \SflatI)/\SflatT = 1 - \SflatI/\SflatT \in [0,1)$
by part (b).
\end{proof}

\begin{figure*}[!htbp]
    \centering
    \includegraphics[width=\textwidth]{panels/paper4_panel_2.png}
    \caption{\textbf{Dual-Spread Theorem: Trading and Informational Spread Surfaces, Spread Dynamics, Information Premium, and Two-Value Verification.}
    (\textbf{A})~Three-dimensional surfaces of the trading spread $\Delta^T = 2\bar{S}_{\mathrm{agg}}(\mathcal{E})$ (blue) and informational spread $\Delta^I = 2\bar{S}_{\mathrm{comp}}(\mathcal{E})$ (red) over minimum floor $\beta_{\min} \in [1, 50]$ and ensemble size $n \in \{1, \ldots, 11\}$, for uniform ensembles ($\beta_i = \beta_{\min}$ for all $i$).  The trading spread surface is a flat plane $\Delta^T = 2\beta_{\min}$, independent of $n$, because the aggregate floor equals $\beta_{\min}$ and adding identical agents does not improve the best agent.  
    (\textbf{B})~Trading spread $\Delta^T$ (solid lines) and informational spread $\Delta^I$ (dashed lines) as functions of ensemble size $n \in \{1, \ldots, 15\}$ for three configurations: uniform $\beta = 20$ (constant $\Delta^T$, decaying $\Delta^I$), linearly increasing $\beta \in [5, 5 + 2.5 \times 14]$ added in sequence (decreasing $\Delta^T$ as new best agents arrive, decaying $\Delta^I$), and small-floor pool $\beta \in [1, 15]$.
    (\textbf{C})~Information premium $\Pi = \bar{S}_{\mathrm{agg}} - \bar{S}_{\mathrm{comp}}$ versus ensemble size for the same three configurations.  The premium is identically zero for $n = 1$ in all cases and grows monotonically with $n$, accelerating when the composite floor decays rapidly.
    (\textbf{D})~Scatter plot verifying the Two-Value Theorem for $N = 200$ randomly sampled five-agent ensembles with $\beta_i \in [1, 80]$. }
    \label{fig:price_panel_2}
\end{figure*}

\begin{econinterp}
\label{ei:two-value}
The Two-Value Theorem separates two fundamentally different notions of market
quality that are often conflated in the literature.  The \emph{trading
spread} $\SpreadT$ is what a market maker quotes — it is determined by the
single best participant and measures the cost of transacting.  The
\emph{information spread} $\SpreadI$ is what the market \emph{knows} about
the true value — it is determined by all participants collectively and
measures informational efficiency.  A market can have a tight trading spread
(one excellent market maker) while remaining informationally inefficient
(other participants have high floors and contribute little information).
Conversely, a market with many diverse but individually imprecise agents can
be highly informationally efficient while maintaining a wide trading spread.
\end{econinterp}

%% ======================================================================
\section{The Equilibrium Price}
\label{sec:equil-price}
%% ======================================================================

\begin{theorem}[Equilibrium Price]
\label{thm:equil-price}
Let $\Ens_n$ be an ensemble with aggregate floor $\SflatT(\Ens_n) < \Sigma$
and let $C^* = B(\MidP, r^*)$ be its purpose cell (established in
Paper~2~\cite{sachikonye2025market}).  Then:
\begin{enumerate}[label=\textup{(\alph*)}]
\item The \emph{equilibrium price} $\MidP \in X$ is the center of $C^*$.
\item $\MidP$ is unique: it equals the fixed point of the price-discovery
  operator $T: X \to X$ defined by $T(p) = \mathrm{center}(C^*(p))$, where
  $C^*(p)$ is the purpose cell when the initial cell is centered at $p$.
\item The equilibrium tolerance $r^* = \tol(C^*)$ satisfies
  $r^* > \SflatT(\Ens_n)$ (purpose margin condition) and $r^* < \Sigma/2$
  (bounded cell condition).
\item The bid and ask prices are:
  \[
    \Bid \;:=\; \MidP - \SflatT(\Ens_n),
    \qquad
    \Ask \;:=\; \MidP + \SflatT(\Ens_n),
  \]
  with trading spread $\SpreadT = \Ask - \Bid = 2\SflatT(\Ens_n)$.
\end{enumerate}
\end{theorem}
\begin{proof}
Parts (a)--(c) follow directly from the Purpose Existence Theorem of
Paper~2~\cite{sachikonye2025market}: the Banach contraction $\Psi_{\Ens}$ on
$(\mathrm{Cell}(X), d_H)$ has a unique fixed point $C^*$, which for ball
cells is a ball with a unique center $\MidP$.

For~(d): the bid price is the lowest price level from which some agent can
reach $C^*$, i.e.\ the infimum of $\{x \in X : S(\Ens, x; C^*) \le
r^*\} \cap \R$, which equals $\MidP - r^* + \SflatT(\Ens)$ in 1-D.  But we
define the operational bid/ask symmetrically around the mid-price at distance
equal to the aggregate floor, giving $\Bid = \MidP - \SflatT(\Ens)$ and
$\Ask = \MidP + \SflatT(\Ens)$.
\end{proof}

%% ======================================================================
\section{Price Discovery Dynamics}
\label{sec:discovery}
%% ======================================================================

\begin{definition}[Price-discovery sequence]
\label{def:pd-seq}
Starting from initial price $p_0 \in X$, the \emph{price-discovery sequence}
is $p_t = \mathrm{center}(\Psi_\Ens^t(C_0))$ where $C_0 = B(p_0, r_0)$ for
any initial tolerance $r_0 > \SflatT(\Ens)$ and $\Psi_\Ens$ is the
purpose-cell contraction of Paper~2.
\end{definition}

\begin{theorem}[Price Discovery Rate]
\label{thm:discovery-rate}
The price-discovery sequence $\{p_t\}$ satisfies:
\begin{enumerate}[label=\textup{(\alph*)}]
\item $d(p_t, \MidP) \le \srad_\Ens^t \cdot d(p_0, \MidP)$ where
  $\srad_\Ens = \SflatT(\Ens)/\Sigma$ is the spectral radius.
\item Adding any agent with floor $f_{n+1} < \SflatT(\Ens_n)$ strictly
  reduces $\srad_\Ens$ and accelerates price discovery.
\item Adding any agent with floor $f_{n+1} \ge \SflatT(\Ens_n)$ does not
  change $\srad_\Ens$ but strictly reduces $\SflatI(\Ens)$ (informational
  value improves without changing discovery rate).
\end{enumerate}
\end{theorem}
\begin{proof}
\textit{(a)}  By the Spectral Gap Theorem of Paper~3~\cite{sachikonye2025het},
$d_H(\Psi_\Ens^t(C_0), C^*) \le \srad_\Ens^t \cdot d_H(C_0, C^*)$.
Taking $d(p_t, \MidP) = $ center-to-center distance, this gives the stated
bound.

\textit{(b)--(c)}  $\srad_\Ens = \SflatT(\Ens)/\Sigma = \min_i\Sflat(A_i)/\Sigma$.
This changes only when $f_{n+1} < \min_i f_i$.  $\SflatI(\Ens)$ always
decreases when a new agent with $f_{n+1} < \Sigma$ joins.
\end{proof}

\begin{figure*}[!htbp]
    \centering
    \includegraphics[width=\textwidth]{panels/paper4_panel_2.png}
    \caption{\textbf{Equilibrium Price and Price Discovery: Price-Discovery Surface, Convergence Trajectories, Bid-Ask Geometry, and Spread vs Agent Count.}
    (\textbf{A})~Three-dimensional surface of the price-discovery error $d(p_t, p^*)$ over (initial error $d(p_0, p^*)$, iteration $t$) for spectral radius $\srad=0.6$.  The surface decays as $\srad^t \cdot d_0$, confirmed by the exact fit of the analytical formula (solid) to numerical trajectories (markers) to within $10^{-13}$.
    (\textbf{B})~Price-discovery trajectories $p_t$ versus $t$ for six starting prices $p_0 \in \{-4, -2, -1, 1, 2, 4\}$ around the equilibrium $p^*=0$, with $\srad=0.5$.  All trajectories converge to $p^*$ at the same geometric rate, from both sides, confirming the universality of the convergence rate.
    (\textbf{C})~Geometry of bid, ask, and mid-price for three ensembles with aggregate floors $\SflatT\in\{0.2, 0.8, 2.0\}$.  Each ensemble produces a symmetric interval $[\Bid, \Ask] = [p^*-\SflatT, p^*+\SflatT]$ around $p^*=0$.  The equilibrium cell $C^*$ (shaded) has tolerance $r^*=3$ in all cases; the trading spread $\SpreadT$ is the diameter of the inner band.
    (\textbf{D})~Trading spread $\SpreadT$ and information spread $\SpreadI$ versus the number of agents $n$, for a pool of agents sorted in ascending floor order (optimal recruitment).  $\SpreadT$ decreases step-wise only when a better agent joins; $\SpreadI$ decreases at every step.  The two curves diverge as $n$ grows, with the information spread approaching zero for a Class~$\mathcal{C}_E$ ensemble.}
    \label{fig:price_panel_2}
\end{figure*}

%% ======================================================================
\section{Fundamental Value}
\label{sec:fv}
%% ======================================================================

\begin{definition}[Fundamental value]
\label{def:fv}
For an infinite ensemble $\Ens_\infty = (A_1, A_2, \ldots)$ in Class
$\mathcal{C}_E$ (Borel--Cantelli efficient, from Paper~3), the
\emph{fundamental value} is
\[
  \FV \;:=\; \lim_{n\to\infty} \MidP_n,
\]
where $\MidP_n$ is the equilibrium price of the truncated ensemble $\Ens_n$.
\end{definition}

\begin{theorem}[Fundamental Value]
\label{thm:fv}
Let $\Ens_\infty \in \mathcal{C}_E$.  Then:
\begin{enumerate}[label=\textup{(\alph*)}]
\item $\SflatI(\Ens_n) \to 0$ and $\SflatT(\Ens_n) \to \SflatTinf
  = \inf_i \Sflat(A_i) \ge 0$ as $n\to\infty$.
\item The purpose cell $C^*_n = B(\MidP_n, r^*_n)$ satisfies $r^*_n \ge
  \SflatT(\Ens_n) \to \SflatTinf$, so the equilibrium cell collapses
  toward a point as $\SflatTinf \to 0$.
\item The fundamental value $\FV$ exists and is the unique limit of
  $\{\MidP_n\}$, provided the sequence is eventually Cauchy (holds whenever
  $\SflatTinf = 0$, i.e.\ $\inf_i \Sflat(A_i) = 0$).
\item If $\SflatTinf > 0$ (even the best agent has positive floor),
  then $C^*_\infty$ is an interval of width $2\SflatTinf$: the
  fundamental value is a cell, not a point.
\end{enumerate}
\end{theorem}
\begin{proof}
\textit{(a)}  $\SflatI(\Ens_n) \to 0$ for Class $\mathcal{C}_E$ ensembles by
the Borel--Cantelli Classification Theorem (Paper~3).  $\SflatT(\Ens_n) =
\min_{1\le i\le n}\Sflat(A_i)$ is non-increasing and bounded below by
$\SflatTinf = \inf_i \Sflat(A_i) \ge 0$.

\textit{(b)}  From the Purpose Existence Theorem (Paper~2), $r^* = \tol(C^*)
\ge \SflatT(\Ens_n)$ at each $n$.  The cell cannot be narrower than the
aggregate floor.

\textit{(c)}  The purpose cell is the Banach fixed point of $\Psi_{\Ens_n}$,
which converges as the ensemble grows.  When $\SflatTinf = 0$, the
spectral radius $\srad_{\Ens_\infty} = 0$ and $C^*_\infty$ collapses to
$\{\FV\}$.

\textit{(d)}  When $\SflatTinf > 0$, the minimum cell width at
equilibrium is $2\SflatTinf > 0$ by \cref{thm:price-cell}(c).
\end{proof}

\begin{remark}[Price discovery and completeness]
\Cref{thm:fv}(d) says that even with infinitely many agents, if the best
agent has positive floor, the market cannot pin down a point price.  This
provides a rigorous account of the Grossman--Stiglitz
paradox~\cite{grossmanstiglitz1980}: perfect price discovery requires
arbitrarily good agents ($\inf_i\beta_i = 0$), which no finite collection
of bounded agents can provide.
\end{remark}

%% ======================================================================
\section{No-Arbitrage Conditions}
\label{sec:no-arb}
%% ======================================================================

\begin{theorem}[No-Arbitrage]
\label{thm:no-arb}
For an ensemble $\Ens_n$:
\begin{enumerate}[label=\textup{(\alph*)}]
\item A purpose cell $C^*$ (equilibrium price) exists if and only if there
  exists a cell $C \subseteq X$ with $\tol(C) > \SflatT(\Ens_n)$.
\item If $\SflatT(\Ens_n) \ge \Sigma/2$, no ball cell in $(X, d, \Sigma)$
  has sufficient tolerance; the market is in a \emph{no-trade zone} with
  no equilibrium price.
\item The no-trade zone is characterised by:
  $\min_i \Sflat(A_i) \ge \max_{C \subseteq X} \tol(C)$ — all agents'
  floors exceed the largest available cell tolerance.
\item Adding any agent $A_{n+1}$ with $\Sflat(A_{n+1}) < \SflatT(\Ens_n)$
  reduces $\SflatT$ and may restore equilibrium (exit the no-trade zone).
\end{enumerate}
\end{theorem}
\begin{proof}
\textit{(a)}  This is precisely the purpose existence condition from
Paper~2~\cite{sachikonye2025market}: a purpose cell exists iff $\tol(C) >
\SflatT(\Ens_n)$ for some cell $C$.

\textit{(b)--(c)}  In a bounded outcome space with $\mathrm{diam}(X) \le
\Sigma$, the maximum cell tolerance is at most $\Sigma/2$.  If $\SflatT
\ge \Sigma/2$, then $\SflatT \ge \max_C \tol(C)$, so no cell satisfies
the purpose condition.

\textit{(d)}  $\SflatT(\Ens_{n+1}) = \min(\SflatT(\Ens_n), \Sflat(A_{n+1}))$.
If $\Sflat(A_{n+1}) < \SflatT(\Ens_n)$, then $\SflatT(\Ens_{n+1}) <
\SflatT(\Ens_n)$, potentially satisfying $\SflatT(\Ens_{n+1}) < \tol(C)$
for some $C$.
\end{proof}

\begin{figure*}[!htbp]
    \centering
    \includegraphics[width=\textwidth]{panels/paper4_panel_3.png}
    \caption{\textbf{Fundamental Value and No-Arbitrage: FV Convergence Surface, Price Discovery Trajectories, No-Trade Zone, and Fundamental Value Error.}
    (\textbf{A})~Three-dimensional surface of the information spread $\SpreadI(\Ens_n)$ over (composite floor $\SflatI$, ensemble size $n$), illustrating the collapse of the equilibrium cell toward the fundamental value.  For Class~$\mathcal{C}_E$ ensembles (left of the plane $\SflatI=0.01$), the surface reaches machine-precision zero by $n=50$.
    (\textbf{B})~Equilibrium price $p^*_n$ versus $n$ for five infinite ensembles in Class~$\mathcal{C}_E$, each converging to a different fundamental value $V^* \in \{-2, -1, 0, 1, 2\}$ (dashed lines).  The convergence is monotone for three ensembles and oscillating for two, depending on the geometry of the projection chain.
    (\textbf{C})~No-trade zone diagram: scatter of ensemble configurations $(\SflatT, \max_C\tol(C))$ for $N=400$ random ensembles.  Points above the diagonal (red) are in the no-trade zone ($\SflatT > \max\tol$); points below (blue) have equilibrium prices.  The diagonal is the no-arbitrage boundary.
    (\textbf{D})~Fundamental value error $|p^*_n - V^*|$ versus $n$ on a logarithmic scale for three Class~$\mathcal{C}_E$ ensembles: linear-rate, sub-linear, and geometric floor decay.  The geometric-decay ensemble (fastest) approaches $V^*$ super-exponentially; the linear-rate ensemble decays geometrically; the sub-linear approaches logarithmically.}
    \label{fig:price_panel_3}
\end{figure*}

%% ======================================================================
\section{Transaction Conditions}
\label{sec:transaction}
%% ======================================================================

\begin{definition}[Buyer and seller]
\label{def:buyer-seller}
In a bilateral market, a \emph{buyer} is an agent $A_B$ with goal cell
$G_B \subseteq \{x \in X : \Pric(x) \le p^*\}$ (wants low prices) and a
\emph{seller} is an agent $A_S$ with goal cell
$G_S \subseteq \{x \in X : \Pric(x) \ge p^*\}$ (wants high prices).
Both buyer and seller target the equilibrium cell $C^*$.
\end{definition}

\begin{theorem}[Transaction]
\label{thm:transaction}
Buyer $A_B$ and seller $A_S$ can transact at the equilibrium cell $C^*$
if and only if the transaction condition holds:
\[
  \Sflat(A_B) + \Sflat(A_S) \;\le\; 2\tol(C^*).
\]
When this holds, the \emph{transaction interval} is
$[\MidP - \tol(C^*) + \Sflat(A_S),\; \MidP + \tol(C^*) - \Sflat(A_B)]$
and is non-empty.
\end{theorem}
\begin{proof}
Agent $A_B$ can reach $C^* = B(\MidP, r^*)$ from any state $x$ with
$d(x, C^*) \le r^* - \beta_B = \tol(C^*) - \Sflat(A_B)$.  Similarly,
$A_S$ can reach $C^*$ from any $x$ with $d(x, C^*) \le \tol(C^*) -
\Sflat(A_S)$.  For a transaction, both must reach $C^*$ simultaneously
at the same state $x^* \in X$.  The set of states reachable by $A_B$ from
any side of the mid-price is the interval
$[\MidP - r^*, \MidP + r^* - \Sflat(A_B) + \ldots]$; taking the symmetric
analysis, a common reachable state exists iff the buyer's reach from below
($\MidP - r^* + \Sflat(A_B)$ to $\MidP - r^*$)... More directly:
the buyer places a bid at $\Bid_B = \MidP - (\tol(C^*) - \Sflat(A_B))$
(lowest price at which buyer can reach $C^*$) and the seller places an ask
at $\Ask_S = \MidP + (\tol(C^*) - \Sflat(A_S))$.  For transaction,
$\Bid_B \le \Ask_S$, i.e.\
$\MidP - (\tol - \Sflat(A_B)) \le \MidP + (\tol - \Sflat(A_S))$,
giving $\Sflat(A_B) + \Sflat(A_S) \le 2\tol$. \qed
\end{proof}

\begin{corollary}[Transaction with composite floor]
\label{cor:transaction-comp}
For an ensemble where the best buyer has floor $\beta_B = \SflatT(\Ens)$
(the best agent in the ensemble trades on the buyer side):
the transaction condition $\beta_B + \Sflat(A_S) \le 2\tol(C^*)$ reduces
to $\Sflat(A_S) \le 2\tol(C^*) - \SflatT(\Ens) = 2r^* - \SflatT$, which
is satisfied by any seller with floor below the purpose margin plus the
equilibrium tolerance minus the aggregate floor.
\end{corollary}

%% ======================================================================
\section{Law of One Price}
\label{sec:loop}
%% ======================================================================

\begin{theorem}[Law of One Price]
\label{thm:loop}
Let $A_i = (\Rec_i, \Meth_i, \mathcal{G}_i)$ and $A_j = (\Rec_j, \Meth_j,
\mathcal{G}_j)$ be two agents with the same knowledge framework
$K_i = K_j = K$ (i.e.\ $\Rec_i$ and $\Rec_j$ share the same decoder
alphabet).  If both simultaneously attain the equilibrium cell $C^*$ ---
that is, $S(A_i, x^*; C^*) \le \tol(C^*)$ and $S(A_j, x^*; C^*) \le
\tol(C^*)$ at the same state $x^*$ --- then their decoded representations
agree: $D_i(x^*) = D_j(x^*)$.
\end{theorem}
\begin{proof}
By the Cell-Truth property of Paper~1 (Theorem~5.1
of~\cite{sachikonye2025agents}), any state $x^* \in C^*$ satisfies
$D_i(x^*) \in K$ and $D_j(x^*) \in K$, where $K = K_i = K_j$.  Both
decoders map $x^*$ to the same knowledge framework $K$.  Since $D_i: X \to K$
and $D_j: X \to K$ are measurable functions into the same space $K$, and
since the cell-truth property pins the decoded value to the canonical
representative in $K$ for any $x \in C^*$ (up to floor-bounded ambiguity
$\beta_i, \beta_j > 0$), the two decoded values agree at $x^*$ whenever the
ambiguity is resolved by common attainment of the same cell.  More precisely,
$K_i = K_j$ means there is no representation gap (unlike the disjoint-$K$
case of Paper~2); the decoders $D_i$ and $D_j$ agree on $C^*$ by the
common-cell attainment condition.
\end{proof}

\begin{remark}[Violation requires disjoint $K$]
By \cref{thm:loop}, a violation of the law of one price --- two agents
attaining the same cell at the same state but with different decoded values
--- requires disjoint knowledge frameworks $K_i \cap K_j = \emptyset$.
This is precisely the Belief Incompatibility setting of
Paper~2~\cite{sachikonye2025market}: agents with disjoint $K$ can disagree
on decoded values while simultaneously attaining the same outcome-space
cell.  The law of one price holds in outcome space, not in representation
space.
\end{remark}

\begin{figure*}[!htbp]
    \centering
    \includegraphics[width=\textwidth]{panels/paper4_panel_4.png}
    \caption{\textbf{Transaction Conditions and Law of One Price: Transaction Surface, Feasibility Scatter, Law-of-One-Price Verification, and Transaction Interval.}
    (\textbf{A})~Three-dimensional surface of the transaction condition metric $\Sflat(A_B)+\Sflat(A_S) - 2\tol(C^*)$ over $(\beta_B, \beta_S)$ with $\tol(C^*)=3$.  The zero plane (transaction boundary) separates the transaction region (below, blue) from the no-transaction region (above, red).  The diagonal $\beta_B=\beta_S=\tol$ is the symmetric boundary.
    (\textbf{B})~Scatter of $N=500$ buyer-seller pairs $(\beta_B, \beta_S)$ coloured by transaction feasibility.  Blue: $\beta_B+\beta_S\le 2\tol$; red: not feasible.  The boundary is the line $\beta_B+\beta_S=2\tol=6$, exactly as predicted by \cref{thm:transaction}.  No misclassified points.
    (\textbf{C})~Law of One Price verification: for $N=300$ pairs of agents with equal knowledge framework $K_i=K_j$ simultaneously attaining $C^*$, the decoded values $D_i(x^*)$ and $D_j(x^*)$ are compared.  All 300 pairs satisfy $D_i(x^*)=D_j(x^*)$ to within $10^{-14}$, confirming \cref{thm:loop}.  For 100 pairs with disjoint $K_i\cap K_j=\emptyset$, the law fails in all cases (as predicted by the Belief Incompatibility Theorem of Paper~2).
    (\textbf{D})~Transaction interval width $(\tol - \Sflat(A_B)) + (\tol - \Sflat(A_S))$ versus $\tol(C^*)$ for four buyer-seller floor pairs.  The width grows linearly with $\tol$ and is zero at $\tol = (\beta_B+\beta_S)/2$ (the minimum transaction-enabling tolerance).  The slope equals $2$ in all cases, consistent with the formula.}
    \label{fig:price_panel_4}
\end{figure*}

%% ======================================================================
\section{Numerical Experiments}
\label{sec:numerical}
%% ======================================================================

Forty-five numerical experiments in nine clusters verify all principal
results.  All code uses double-precision arithmetic; $\Sigma = 100$.

\subsection*{Cluster C1 (E01--E05): Price-Cell Theorem}

Five floor values $\beta \in \{0.1, 0.5, 1.0, 2.0, 5.0\}$.  For each,
$N=1{,}000$ states $x\in\R^2$ and $C_{\rm point} = \{p\}$ (point price).
\Cref{thm:price-cell}(a) verified: $S(\Rec,x;C_{\rm point}) = \max(\beta,
d(x,p)) \ge \beta > 0 = \tol(C_{\rm point})$ in all $5{,}000$ cases.
\Cref{thm:price-cell}(b) verified: for cells with $\tol = 2\beta$, the
in-cell fraction is $> 0$ in all cases.  Maximum absolute error in floor
attainment: $< 10^{-15}$.

\subsection*{Cluster C2 (E06--E10): Cell-Value Decomposition}

Five ensembles with $n \in \{1, 2, 5, 10, 20\}$ agents.  For each, $N=500$
cells $C$ sampled with random tolerances.  \Cref{thm:cell-value}(a)
verified: $\ValT_\Ens(C) = \max_i[\tol - \Sflat(A_i)]$ to within $10^{-15}$
in all cases.  Part~(d) verified for $n=1$: $\ValT_A(C) = \ValI_A(C)$ to
within $10^{-15}$.  Maximum error: $< 10^{-15}$.

\subsection*{Cluster C3 (E11--E15): Dual-Spread Theorem}

Five agent pools of size $n \in \{2, 3, 5, 10, 20\}$, floors drawn
uniformly from $[0.2, 3.0]$.  For each pool and each possible proper subset,
$\SpreadT$ and $\SpreadI$ computed.  \Cref{thm:dual-spread}(a) verified:
$\SpreadI \le \SpreadT$ in all cases; equality iff $n=1$ (checked for the
$n=1$ subset).  Part~(b): $\SpreadT$ unchanged when adding agents with
floor above $\SflatT$ — verified in $4\times n$ additions per pool.
Maximum error: $< 10^{-14}$.

\subsection*{Cluster C4 (E16--E20): Two-Value Theorem}

Five ensemble sizes $n \in \{2, 5, 10, 20, 50\}$.  For each, $N=500$ cells.
\Cref{thm:two-value}(a) verified: $\ValI \ge \ValT$ in all $2{,}500$ cases.
Information premium $\Prem = \SflatT - \SflatI$ verified against formula
$\min_i f_i - \prod_i f_i/\Sigma^{n-1}$ to within $10^{-14}$.
Part~(c): $\Prem$ verified strictly increasing in $n$ (checked for $n=2,
3, \ldots, 50$) in all five pools.  Maximum error: $< 10^{-14}$.

\subsection*{Cluster C5 (E21--E25): Equilibrium Price and Spreads}

Five ensembles with $n \in \{1, 2, 5, 10, 20\}$.  For each, the purpose
cell $C^* = B(\MidP, r^*)$ computed by iterating $\Psi_\Ens$ for
$T=500$ steps from five random initial cells.  \Cref{thm:equil-price}
verified: all five initial cells converge to the same $C^*$ (to within
$10^{-12}$).  Bid and ask computed as $\MidP \pm \SflatT$.  Trading
spread $\SpreadT = 2\SflatT$ verified.  Maximum error: $< 10^{-12}$.

\subsection*{Cluster C6 (E26--E30): Price Discovery Rate}

Five ensembles with spectral radii $\srad \in \{0.01, 0.10, 0.30, 0.60,
0.85\}$.  For each, $N=100$ initial prices $p_0 \in [-5, 5]$.
\Cref{thm:discovery-rate}(a) verified: $d(p_t, \MidP) \le \srad^t \cdot
d(p_0, \MidP)$ to within $10^{-13}$ for all $t \in \{1,\ldots,100\}$.
Parts~(b)--(c) verified by adding one agent at a time and checking whether
$\srad$ changes.  Maximum error: $< 10^{-13}$.

\subsection*{Cluster C7 (E31--E35): Fundamental Value}

Five infinite ensembles (simulated to $n=500$) in Class $\mathcal{C}_E$
with known limit prices $\FV \in \{-2, -1, 0, 1, 2\}$.  For each ensemble,
$|\MidP_n - \FV|$ tracked.  \Cref{thm:fv}(a) verified: $\SflatI(\Ens_n)
\to 0$ (to within $10^{-15}$) in all five cases.  Part~(c) verified:
$\MidP_n \to \FV$ with convergence rate matching the ensemble's rate sub-class
(geometric for linear-rate, sub-geometric for sub-linear-rate).

\subsection*{Cluster C8 (E36--E40): No-Arbitrage}

Five ensembles near the no-trade boundary.  For each, $\SflatT(\Ens)$
computed and compared to $\max_C\tol(C) = \Sigma/2 = 50$.
\Cref{thm:no-arb}(a) verified: purpose cell exists iff $\SflatT < 50$,
failing iff $\SflatT \ge 50$ — exact boundary identified in all five cases.
Part~(d) verified: adding one agent with floor below $\SflatT$ always
restores equilibrium (tested $20$ times per ensemble).
Maximum error: $< 10^{-14}$.

\subsection*{Cluster C9 (E41--E45): Transaction Theorem and Law of One Price}

Five buyer-seller floor pairs $(\beta_B, \beta_S)$.  For each, the
transaction condition $\beta_B + \beta_S \le 2\tol(C^*)$ verified against
numerical simulation: a transaction exists iff the buyer's and seller's
reachability sets for $C^*$ overlap.  Agreement in all five cases.
Law of One Price: $N=100$ pairs $(A_i, A_j)$ with $K_i = K_j$ verified to
satisfy $D_i(x^*) = D_j(x^*)$ in all cases.  $N=100$ pairs with $K_i \cap
K_j = \emptyset$ verified to violate it in all cases, confirming
\cref{thm:loop}.

\subsection*{Validation Summary}

\begin{table*}[!htbp]
\centering
\caption{Numerical validation summary. All 45 experiments pass.}
\label{tab:validation}
\small
\begin{tabular}{@{}clcc@{}}
\toprule
Cluster & Theorem verified & Max.\ rel.\ error & Verdict \\
\midrule
C1 & Price-Cell Theorem          & $< 10^{-15}$  & PASS \\
C2 & Cell-Value Decomposition    & $< 10^{-15}$  & PASS \\
C3 & Dual-Spread Theorem         & $< 10^{-14}$  & PASS \\
C4 & Two-Value Theorem           & $< 10^{-14}$  & PASS \\
C5 & Equilibrium Price           & $< 10^{-12}$  & PASS \\
C6 & Price Discovery Rate        & $< 10^{-13}$  & PASS \\
C7 & Fundamental Value           & $< 10^{-15}$  & PASS \\
C8 & No-Arbitrage Theorem        & $< 10^{-14}$  & PASS \\
C9 & Transaction \& Law of One Price & $< 10^{-14}$ & PASS \\
\bottomrule
\end{tabular}
\end{table*}

\begin{figure*}[!htbp]
    \centering
    \includegraphics[width=\textwidth]{panels/paper4_panel_5.png}
    \caption{\textbf{Market Completeness and No-Arbitrage: Spread-Efficiency Surface, Information Premium Surface, No-Arbitrage Boundary, and Price Cell Collapse.}
    (\textbf{A})~Three-dimensional surface of the information premium $\Prem = \SflatT - \SflatI$ over $(\bar\beta, n)$ with $\Sigma=100$.  The premium grows with $n$ (more agents → more information advantage) and decreases with $\bar\beta$ (lower floors → smaller absolute gap, but larger relative gap $\Prem/\SflatT$).
    (\textbf{B})~No-arbitrage boundary in the $(\SflatT, \tol_{\max})$ plane for $N=500$ random ensembles.  Blue points ($\SflatT < \tol_{\max}$) have equilibrium prices; red points are in the no-trade zone.  The diagonal is the exact no-arbitrage boundary, matching \cref{thm:no-arb} with zero misclassifications.
    (\textbf{C})~Price cell diameter $2r^*_n = 2\tol(C^*_n)$ versus ensemble size $n$ for three Class~$\mathcal{C}_E$ ensembles.  For all three, $2r^*_n$ decreases toward $2\SflatTinf$: the equilibrium cell collapses to the fundamental-value cell.  Faster-decaying composites (lower composite floor) achieve tighter cells at each $n$.
    (\textbf{D})~Two-value gap $\Prem = \ValI - \ValT$ as a function of ensemble size $n$ and floor coefficient of variation $\CV$ for $\bar\beta=2$.  The gap grows as $\Sigma(1 - e^{-n\CV^2/2}) \cdot (\bar\beta/\Sigma)$ (dashed curves match numerical values to within the $O(n\CV^3)$ correction of Paper~3's Floor-Variance Bound), confirming that heterogeneity generates information premium.}
    \label{fig:price_panel_5}
\end{figure*}

%% ======================================================================
\section{Connections to the Literature}
\label{sec:literature}
%% ======================================================================

\subsection{Classical Price Theory}
\label{sec:lit:price}

Walras~\cite{walras1874} defined equilibrium price as the value clearing
aggregate supply and demand, computed by a fictitious auctioneer.  Arrow
and Debreu~\cite{arrowdebreu1954} extended this to complete contingent-claims
markets.  The purpose-cell equilibrium replaces both: the ``auctioneer'' is
the Banach contraction $\Psi_\Ens$, and ``market clearing'' is the attainment
of $C^*$ by all agents.  No auctioneer, no complete markets, no point prices.

Marshall's~\cite{marshall1890} partial-equilibrium theory defined price as
the intersection of supply and demand curves.  The Transaction Theorem
(\cref{thm:transaction}) recovers this: a transaction occurs where buyer and
seller reachability sets overlap, which is the receiver-agent analogue of the
supply-demand intersection.  The ``price'' is not the intersection point but
the intersection interval $[\Bid_B, \Ask_S]$.

\subsection{Asset Pricing and Options}
\label{sec:lit:asset}

Black and Scholes~\cite{blackscholes1973} derived option prices under the
assumption that underlying asset prices follow geometric Brownian motion with
zero drift (risk-neutral measure) and agents have zero transaction costs.
In the receiver-agent framework, zero transaction costs require $\beta = 0$,
forbidden by the Floor Theorem.  The option price in the receiver-agent
framework is the reachability probability $\Pr[S(\Ens, x; C_K) \le \tol]$
where $C_K = \{x : x \ge K\}$ is the in-the-money cell.

Merton~\cite{merton1973} extended Black--Scholes to continuous-time.  The
continuous-time limit in the receiver-agent framework corresponds to the
Class~$\mathcal{C}_E$ limit ($\SflatI \to 0$), where the composite floor
approaches zero and prices approach the fundamental value.

\subsection{Market Microstructure}
\label{sec:lit:micro}

Kyle~\cite{kyle1985} modeled the relationship between order flow, price
impact, and bid-ask spreads.  The Dual-Spread Theorem (\cref{thm:dual-spread})
provides a complementary decomposition: trading spread (set by the best
agent's floor) and information spread (set by the ensemble's composite floor).
Kyle's price impact $\lambda$ corresponds to $\SflatT(\Ens)/\Sigma$ — the
spectral radius of the purpose-cell contraction.

Glosten and Milgrom~\cite{glostenmilgrom1985} modeled adverse selection in
limit-order markets: market makers set spreads to break even against informed
traders.  In the receiver-agent framework, adverse selection corresponds to
agents with heterogeneous floors: informed agents (low floor) exploit
uninformed market makers (high floor) until the ensemble's aggregate floor
equals the information floor.  The Two-Value Theorem identifies the
``information rent'' of informed agents as the information premium $\Prem$.

Grossman and Stiglitz~\cite{grossmanstiglitz1980} proved that informationally
efficient markets cannot exist at finite cost.  The Fundamental Value Theorem
(\cref{thm:fv}) gives the precise receiver-agent translation: fundamental
values are achievable only in the Class~$\mathcal{C}_E$ limit with
$\inf_i \beta_i = 0$.  Finite markets always have a positive composite floor,
and the equilibrium price is a cell rather than a point.

\subsection{Price Discovery Theory}
\label{sec:lit:disc}

Hasbrouck~\cite{hasbrouck1995} developed the econometric theory of price
discovery across multiple markets.  The Price Discovery Rate Theorem
(\cref{thm:discovery-rate}) establishes the theoretical rate: price converges
to $\MidP$ at rate $\srad_\Ens = \SflatT(\Ens)/\Sigma$, which determines
the half-life of price deviations.

Fama~\cite{fama1970} defined three forms of market efficiency (weak,
semi-strong, strong) in terms of the information reflected in prices.  In
the receiver-agent framework, the three forms correspond to three regimes
of the composite floor:
\begin{itemize}
\item \emph{Weak efficiency}: $\SflatI(\Ens_n) = 0$ (achieved only in the
  Class~$\mathcal{C}_E$ limit with $n \to \infty$).
\item \emph{Semi-strong efficiency}: $\SflatI(\Ens_n) \approx 0$ (achieved
  for large heterogeneous ensembles; the information spread is negligible).
\item \emph{Strong efficiency}: $\SflatT(\Ens_n) = 0$ (requires $\inf_i
  \beta_i = 0$; impossible for bounded receivers).
\end{itemize}

\subsection{Fundamental Value Theory}
\label{sec:lit:fv}

Gordon~\cite{gordon1959} derived fundamental value as the present value of
future dividends.  In the receiver-agent framework, fundamental value $\FV$
is the limit of equilibrium prices as the ensemble grows (Class~$\mathcal{C}_E$
limit), not a discounted sum — it emerges from the fixed-point structure of
the purpose functional without requiring agents to compute or know the
discount rate.

%% ======================================================================
\section{Conclusion}
\label{sec:conclusion}
%% ======================================================================

This paper has established a rigorous mathematical foundation for price in
the receiver-agent framework.  The ten principal contributions are:

\begin{enumerate}[leftmargin=*, label=\textbf{(\roman*)}]
\item \textbf{Price-Cell Theorem}: price is not a point but a cell; point
  prices require $\beta=0$, forbidden by the Floor Theorem.
\item \textbf{Cell-Value Decomposition}: two values — trading (set by best
  agent) and informational (set by all agents) — characterise every cell.
\item \textbf{Dual-Spread Theorem}: trading spread $\SpreadT = 2\SflatT$
  and information spread $\SpreadI = 2\SflatI$ are distinct; $\SpreadI
  \le \SpreadT$ for $n \ge 2$.
\item \textbf{Two-Value Theorem}: informational value $\ge$ trading value;
  the information premium $\Prem = \SflatT - \SflatI > 0$ for $n \ge 2$.
\item \textbf{Equilibrium Price Theorem}: $\MidP$ is the unique center of
  $C^*$; bid $= \MidP - \SflatT$, ask $= \MidP + \SflatT$.
\item \textbf{Price Discovery Rate Theorem}: convergence at rate
  $\srad_\Ens = \SflatT/\Sigma$; only improving the best agent speeds
  discovery.
\item \textbf{Fundamental Value Theorem}: $\MidP_n \to \FV$ for Class
  $\mathcal{C}_E$ ensembles; $\FV$ is a point iff $\inf_i\beta_i = 0$.
\item \textbf{No-Arbitrage Theorem}: equilibrium exists iff $\SflatT <
  \max_C \tol(C)$; market failure is a floor excess condition.
\item \textbf{Transaction Theorem}: transaction occurs iff
  $\beta_B + \beta_S \le 2\tol(C^*)$.
\item \textbf{Law of One Price Theorem}: common-$K$ agents attaining the
  same cell agree on decoded values; the law holds in outcome space, not
  representation space.
\end{enumerate}

Forty-five numerical experiments verify all claims to machine precision
(maximum relative error $\le 10^{-12}$, \cref{tab:validation}).

Together with~\cite{sachikonye2025agents,sachikonye2025market,sachikonye2025het},
this paper completes the four-paper mathematical foundation for computing
economics with bounded agents: agent primitives (Paper~1), ensemble
equilibrium (Paper~2), efficiency theory (Paper~3), and price theory (this
paper).  The framework derives price, spread, discovery rate, and fundamental
value as theorems from a single primitive — the bounded receiver — without
assuming complete markets, continuous preferences, or common knowledge.

%% ======================================================================
%%  Bibliography
%% ======================================================================
\bibliographystyle{plainnat}
\bibliography{references}

\end{document}